% ============================================================
% ch19.tex —— 第 19 章 电磁对偶（Electric-magnetic duality）
% 底本：Springer LNM 2386，印刷 pp.269–286（PDF p309–p326）
% 编号公式：19.1–19.21（共 21 条，逐条对 200dpi 图像核对）
% ============================================================

\chapter{电磁对偶}

% [原文 p.269]

\ChapterAbstract{我们研究 Maxwell 理论中的电磁对偶，从 theta 角 $\Theta$ 为零的情形开始，再推广到任意值。单极子与双子可以作为 Georgi--Glashow 模型的孤子构型构造出来，这就是 't Hooft--Polyakov 与 Julia--Zee 的构造。它们的质量被 Bogomol'nyi 界从下方界定，而在 Bogomol'nyi--Prasad--Sommerfield（BPS）极限下该界可以饱和，由此产生 BPS 态的概念。这些想法自然地推广到非交换 Yang--Mills 理论，导致 Montonen--Olive 对偶及其与 Langlands 对偶李群的联系。然而，纯 Yang--Mills 理论的强耦合区域仍然难以触及，而重整化又进一步使对偶的量子表述复杂化。$\mathcal{N}=2$ 超对称赋予 BPS 态非凡的稳定性，从而大大深化了电磁对偶的研究，为强耦合规范动力学开了一扇窗。在四维中，$\mathcal{N}=4$ 超 Yang--Mills（SYM）理论既是超共形的又是标度不变的，这些特性使作为猜想的 Montonen--Olive 对偶——在该语境中称为 S 对偶——显得特别自然。把 S 对偶推广到 $\mathcal{N}=2$ 理论（它涵盖更广、更精细的一类量子场论）已带来重要的洞见，尽管其技术难度远比 $\mathcal{N}=4$ 情形大。}

\section{无 $\Theta$ 项的 Maxwell 理论}

\subsection{电磁对偶}

无 $\Theta$ 项的 Maxwell 理论的作用量为：
\begin{equation}
S_M[A;e]=-\frac{1}{4e^2}\int \dd^4x\,F_{\mu\nu}F^{\mu\nu}
=-\frac{1}{2e^2}\int F\wedge\star F . \tag{19.1}
\end{equation}

% [原文 p.270]

电磁张量 $F=\dd A$ 是时空上主 $\mathrm{U}(1)$ 丛上 $\mathrm{U}(1)$-联络 $A$ 的曲率。它满足 Bianchi 恒等式 $\dd F=0$，这解释了 Maxwell 方程组的前一半。后一半由运动方程构成，即 $\dd\star F=0$，其中 $\star$ 是 Hodge 星算子。当存在荷电物质、产生电 1-形式流 $j_e$ 时，运动方程修改为 $\dd\star F=\star j_e$。用分量写出，$F$ 及其对偶 $\star F$ 为：
\[
F^{\mu\nu}=
\begin{pmatrix}
0 & -E_x & -E_y & -E_z\\
E_x & 0 & -B_z & B_y\\
E_y & B_z & 0 & -B_x\\
E_z & -B_y & B_x & 0
\end{pmatrix},
\qquad
\star F^{\mu\nu}=
\begin{pmatrix}
0 & -B_x & -B_y & -B_z\\
E_x & 0 & E_z & -E_y\\
E_y & -E_z & 0 & E_x\\
E_z & E_y & -B_x & 0
\end{pmatrix},
\]
其中 $\mathbf{E}$ 与 $\mathbf{B}$ 分别是电场与磁场。Maxwell 作用量在形如 $A\to A+\ii g^{-1}\dd g$ 的规范变换下不变，其中 $g$ 是从时空到规范群 $\mathrm{U}(1)=\{\ee^{\ii\alpha}\mid \alpha\sim\alpha+2\pi\}$ 的映射。用角 $\alpha$ 表示，规范变换成为 $A\to A-\dd\alpha$。
% [Tutor Note: 原文 p.270 印出的 $\star F^{\mu\nu}$ 矩阵有四个元素与标准的 Hodge 对偶不符：它甚至不是反对称的（$\star F^{01}=-B_x$，而 $\star F^{10}=E_x$）。约定 $F^{0i}=-E_i$、$F^{ij}=-\epsilon^{ijk}B_k$ 时，Hodge 对偶应为 $\star F^{0i}=-B_i$、$\star F^{ij}=\epsilon^{ijk}E_k$，即第二、三、四行的第一列应分别为 $B_x,B_y,B_z$，且第四行第三列应为 $-E_x$。可用"在 Lorentz 号差下 2-形式上 $\star\star=-1$"以及反对称性自查。译文照印原文矩阵。]

当 $j_e=0$ 时，Maxwell 方程在交换 $F$ 与 $\star F$（等价地 $(\mathbf{E},\mathbf{B})\mapsto(\mathbf{B},-\mathbf{E})$）下不变。存在物质时，只要在通常的电 1-形式流 $j_e$ 之外再引入一个磁 1-形式流 $j_m$，并把 Maxwell 方程修改为
\begin{align*}
\dd\star F&=\star j_e,\\
\dd F&=\star j_m,
\end{align*}
这一不变性仍然成立。

\begin{definition}[19.1]
Maxwell 理论的电磁对偶由下述联合变换构成：
\[
F\longmapsto\star F,\qquad (j_e,j_m)\longmapsto(j_m,-j_e).
\]
\end{definition}

点状源的电荷 $q_{\mathrm{el}}=n_e e$（$n_e\in\Z$）可以用 Gauss 定律确定。对任意在固定时刻包围该源的 2-球面 $S^2$，电荷由
\begin{equation}
\int_{S^2}\frac{1}{e}\mathbf{E}\cdot \dd\mathbf{s}
=\frac{1}{e}\int_{S^2}\star F=q_{\mathrm{el}} \tag{19.2}
\end{equation}
给出。

\begin{remark}[19.2]
整数 $n_e$ 标记相应场所在的规范群 $\mathrm{U}(1)$ 的表示。归根结底，电荷的量子化之所以出现，是因为规范群 $\mathrm{U}(1)$ 是紧的。相反，在规范群为非紧群 $\R$ 的类似理论中，电荷未必量子化。
\end{remark}

\begin{definition}[19.3]
磁单极子是磁场 $\mathbf{B}$ 的点状源。与式 (19.2) 中电荷的 Gauss 定律类似，对任意在固定时刻包围磁荷为 $q_{\mathrm{mag}}$ 的磁单极子的 2-球面 $S^2$，

% [原文 p.271]

磁荷可以由下式确定：
\begin{equation}
\int_{S^2}\frac{1}{e}\mathbf{B}\cdot \dd\mathbf{s}
=\frac{1}{e}\int_{S^2}F=q_{\mathrm{mag}} . \tag{19.3}
\end{equation}
\end{definition}

磁荷也是量子化的，由 $q_{\mathrm{mag}}=2\pi n_m/e$ 给出，其中 $n_m\in\Z$。这一量子化来自拓扑：整数 $n_m$ 表示磁单极子所定义的主 $\mathrm{U}(1)$-丛限制到 $S^2$ 上的第一 Chern 类。相反，对规范群 $\R$，$S^2$ 上每个主 $\R$-丛都是平凡的。因此，在规范群为 $\R$ 的理论中没有磁单极子。

\begin{proposition}[19.4]（Dirac 量子化条件）
在 Maxwell 理论中，若存在磁荷为 $q_{\mathrm{mag}}$ 的磁单极子与电荷为 $q_{\mathrm{el}}$ 的带电粒子，则必有
\[
q_{\mathrm{el}}q_{\mathrm{mag}}\in 2\pi\Z .
\]
\end{proposition}

\begin{proof}
这一要求出现的原因之一是：为使电荷为 $q_{\mathrm{el}}$ 的电性源在存在磁荷为 $q_{\mathrm{mag}}$ 的磁单极子时其相位有良好定义（例如见 [346]）。
\end{proof}

与磁单极子类似，在 Maxwell 理论中也可以引入双子：

\begin{definition}[19.5]
双子是同时携带电荷与磁荷的粒子。
\end{definition}

推广命题 19.4，有：

\begin{proposition}[19.6]（Dirac--Schwinger--Zwanziger（DSZ）量子化条件）
对任意两个电荷为 $(q_{\mathrm{el}}^1,q_{\mathrm{mag},1})$ 与 $(q_{\mathrm{el}}^2,q_{\mathrm{mag},2})$ 的双子，必有
\[
q_{\mathrm{el}}^1q_{\mathrm{mag},2}-q_{\mathrm{el}}^2q_{\mathrm{mag},1}\in 2\pi\Z .
\]
\end{proposition}

写 $q_{\mathrm{el}}^i=n_e^i e$ 与 $q_{\mathrm{mag},i}=2\pi n_{m,i}/e$（$i=1,2$），DSZ 量子化条件改写为
\[
n_e^1n_{m,2}-n_e^2n_{m,1}\in\Z .
\]

为在 $F$ 与 $\star F$ 的交换下保持式 (19.2) 与 (19.3) 的归一化，约定
\[
F_D=\frac{2\pi}{e^2}\star F
\qquad\text{与}\qquad
e_D^2=\frac{4\pi^2}{e^2}.
\]

\begin{definition}[19.7]
电磁对偶是这样一个论断：作用量 $S_M[A;e]$ 与 $S_M[A_D;e_D]$ 描述同一物理，其中 $S_M[A;e]$ 由式 (19.1) 给出，而
\[
S_M[A_D;e_D]=-\frac{1}{2e_D^2}\int F_D\wedge\star F_D,
\]
其中 $A_D$ 是曲率为 $F_D$ 的 $\mathrm{U}(1)$ 联络。
\end{definition}

% [原文 p.272]

这一运算一般是对偶而非对称，因为它把耦合为 $e$ 的 Maxwell 理论与耦合为 $e_D$ 的 Maxwell 理论联系起来。注意在电磁对偶下，双子的电荷按 $(n_e,n_m)\longrightarrow(-n_m,n_e)$ 变换；特别地，电子变成单极子、反之亦然。重要的是，DSZ 配对 $n_e^1n_{m,2}-n_e^2n_{m,1}$ 在对偶下保持不变。

\subsection{来自非交换规范理论的单极子与双子}

上一段中我们讨论了 Maxwell 理论中的磁单极子如何可以解释为拓扑非平凡的交换规范场构型。这些单极子自然地作为具有非交换紧规范群的理论中的孤子出现，其中电磁学来自规范对称性的自发破缺。这里我们概述这一构造背后的逻辑，更多细节请参阅 [91, 92, 18, 190, 354, 346]。

考虑带伴随表示中标量场 $\phi$ 的 $\SU(2)$ Yang--Mills 理论，其拉格朗日量定义为
\begin{equation}
\mathcal{L}=-\frac{1}{2g^2}\tr F_{\mu\nu}F^{\mu\nu}
+\frac{1}{g^2}\tr D_\mu\phi D^\mu\phi
-\frac{\lambda}{4}\left(\tr\phi^2-\frac{v^2}{2}\right)^2, \tag{19.4}
\end{equation}
其中 $v$ 是实参数。这个理论称为 Georgi--Glashow 模型。它具有由标量真空期望值 $\langle\phi\rangle$ 参数化的经典真空模空间，满足 $\tr\langle\phi\rangle^2=v^2/2$。写 $\phi=\phi_i\sigma^i/2$，方程即为 $\sum\phi_i^2=v^2$，从而定义 $\R^3$ 中半径为 $v$ 的 2-球面 $S^2_{\mathrm{vac}}$。可以证明任何这样的真空都规范等价于
\[
\langle\phi\rangle=\frac{1}{2}\begin{pmatrix} v & 0\\ 0 & -v\end{pmatrix},
\]
因此规范群总是自发破缺到 $\mathrm{U}(1)$，交换规范耦合为 $e=g$。

考虑有限能量的静态场构型。为使构型具有有限能量，标量势必须在无穷远处足够快地趋于零；特别地，$\phi$ 在无穷远处须满足 $\tr\phi^2\sim v^2/2$。这给出一个映射
\begin{equation}
S^2_\infty\longrightarrow S^2_{\mathrm{vac}}, \tag{19.5}
\end{equation}
其中 $S^2_\infty$ 是渐近空间方向的球面。拓扑上，这样的映射由其指数 $\nu\in\pi_2(S^2)\simeq\Z$ 分类。当指数 $\nu$ 非零时，未破缺子群 $\mathrm{U}(1)\subset\SU(2)$ 依赖于渐近方向；低能交换规范场 $A^{\mathrm{U}(1)}_\mu$ 可以表示为
\begin{equation}
A^{\mathrm{U}(1)}_\mu=\frac{1}{v}\tr(\phi A_\mu), \tag{19.6}
\end{equation}

% [原文 p.273]

其中 $A_\mu$ 是 $\SU(2)$ 规范场。此外，标量 $\phi$ 在无穷远处必须是渐近协变常值的：
\[
D_\mu\phi\sim_\infty 0
\quad\Longrightarrow\quad
A_\mu\sim_\infty \frac{\ii}{v^2}[\phi,\partial_\mu\phi]+\frac{\phi A^{\mathrm{U}(1)}_\mu}{v}.
\]
未破缺交换场 $A^{\mathrm{U}(1)}_\mu$ 的场强 $F^{\mathrm{U}(1)}_{\mu\nu}$ 为
\[
F^{\mathrm{U}(1)}_{\mu\nu}
=\frac{1}{v}\tr(\phi F_{\mu\nu})
\sim_\infty \partial_\mu A^{\mathrm{U}(1)}_\nu-\partial_\nu A^{\mathrm{U}(1)}_\mu
+\frac{\ii}{v^3}\tr(\phi[\partial_\mu\phi,\partial_\nu\phi]),
\]
并满足 $\dd F^{\mathrm{U}(1)}=\dd\star F^{\mathrm{U}(1)}=0$。这一表达式的最后一项来自嵌入 $\mathrm{U}(1)\subset\SU(2)$ 对渐近方向的依赖。可以算得
\[
\int_{S^2_\infty}F^{\mathrm{U}(1)}
=\int_{S^2_\infty}\mathbf{B}^{\mathrm{U}(1)}\cdot \dd\mathbf{s}=4\pi\nu,
\]
其中 $\nu$ 是式 (19.5) 中映射的指数。换言之：

\begin{proposition}[19.8]
Georgi--Glashow 模型中指数 $\nu\ne 0$ 的有限能量静态构型给出未破缺交换 $\mathrm{U}(1)$ 规范群的磁单极子，其磁荷为 $q_{\mathrm{mag}}=2\nu/e$。
\end{proposition}

对这样的孤子可以写出一个拟设，所得单极子称为 't Hooft--Polyakov 单极子 [209, 305]。

\begin{remark}[19.9]
有限能量静态构型的磁数 $n_m$ 是指数 $\nu$ 的两倍，这一事实与我们出发时所用的非交换规范群是 $\SU(2)$ 而非 $\mathrm{SO}(3)$ 有关。$\SU(2)$ 基本表示中的物质会在低能理论中给出电荷为 $1/2$ 的场，相对于这些场，$m=2$ 是 Dirac 量子化条件所允许的最小值。相反，$\mathrm{SO}(3)$ 的表示中的场总是给出整数电荷，与磁荷为 1 的磁单极子的存在相容。我们将在 19.3 节回到这一点。
\end{remark}

在式 (19.4) 定义的 Georgi--Glashow 模型中，还可以构造对低能未破缺 $\mathrm{U}(1)$ 同时携带电荷与磁荷的有限能量静态构型。它们称为 Julia--Zee 双子 [219]。与 't Hooft--Polyakov 单极子的特殊情形一样，未破缺交换规范场的场强写作
\[
F^{\mathrm{U}(1)}_{\mu\nu}=\frac{1}{v}\tr(\phi F_{\mu\nu}).
\]

\subsection{Bogomol'nyi 界与 BPS 极限}

写 $E_i=F_{0i}$、$2B_i=-\epsilon_{ijk}F^{jk}$，则 Georgi--Glashow 模型中 Julia--Zee 双子的质量满足 [219, 59]：

% [原文 p.274]

\begin{equation}
\begin{aligned}
M&=\frac{1}{g^2}\int \dd^3x\left[\tr\left(E_i^2+B_i^2+(D_i\phi)^2+(D_0\phi)^2\right)+V(\phi)\right]\\
&=\frac{1}{g^2}\int \dd^3x\left[\tr\left((E_i-\sin(\theta)D_i\phi)^2+(B_i-\cos(\theta)D_i\phi)^2+(D_0\phi)^2\right)+V(\phi)\right]\\
&\quad+\frac{2}{g^2}\int \dd^3x\left[\sin(\theta)\tr(E_iD^i\phi)+\cos(\theta)\tr(B_iD^i\phi)\right],
\end{aligned} \tag{19.7}
\end{equation}
其中 $\theta$ 是任意角。此外，
\begin{equation}
\frac{1}{g}\int \dd^3x\,\tr(E_iD^i\phi)=vq_{\mathrm{el}},
\qquad
\frac{1}{g}\int \dd^3x\,\tr(B_iD^i\phi)=vq_{\mathrm{mag}}, \tag{19.8}
\end{equation}
其中 $(q_{\mathrm{el}},q_{\mathrm{mag}})$ 是双子的电磁荷。由于 $V(\phi)$ 正定，对使 $\tan\theta=q_{\mathrm{el}}/q_{\mathrm{mag}}$ 的 $\theta$ 可得：

\begin{proposition}[19.10]（Bogomol'nyi 界）
电荷为 $(q_{\mathrm{el}},q_{\mathrm{mag}})$ 的双子的质量 $M_{q_{\mathrm{el}},q_{\mathrm{mag}}}$ 满足界 [59]：
\begin{equation}
M_{q_{\mathrm{el}},q_{\mathrm{mag}}}\ge \frac{2v}{g}\sqrt{q_{\mathrm{el}}^2+q_{\mathrm{mag}}^2}. \tag{19.9}
\end{equation}
\end{proposition}

这一界的一个直接推论是：当 $g=e$ 很小时，磁单极子非常重，而双子更重 [219, 18]。作为对照，低能电子即 $W^\pm$-玻色子的质量通常是 $M_W\sim ve$ 的量级。
% [Tutor Note: 与本书导言（"动机与历史回顾"）中的 Bogomol'nyi 界对照：那里写 $M_{q_{\mathrm{el}},q_{\mathrm{mag}}}\ge\alpha\sqrt{q_{\mathrm{el}}^2+q_{\mathrm{mag}}^2}$，此处式 (19.9) 给出 $\alpha=2v/g$（在 $g=e$ 的 Georgi--Glashow 模型中即 $\alpha=2v/e$），两种写法只差这一常数归一化。另外，(19.7) 第三行中 $\tr(E_iD^i\phi)$ 的指标一上一下即求和约定，可与 (19.8) 的 $vq_{\mathrm{el}}$、$vq_{\mathrm{mag}}$ 对照自查。]

式 (19.7) 中的项 $(D_0\phi)^2$ 与 $V(\phi)$ 表明，Bogomol'nyi 界只有在
\[
D_0\phi=0,\qquad V(\phi)=0
\]
时才能饱和。而这又只有在式 (19.4) 中的耦合 $\lambda$ 恒为零时才可能。

\begin{definition}[19.11]
Georgi--Glashow 模型的这一极限称为 Bogomol'nyi--Prasad--Sommerfield（BPS）极限 [306]。此时理论的拉格朗日量简化为
\[
\mathcal{L}=-\frac{1}{2g^2}\tr F_{\mu\nu}F^{\mu\nu}+\frac{1}{g^2}\tr D_\mu\phi D^\mu\phi,
\]
即理论就是 $\SU(2)$ Yang--Mills 场与一个伴随表示中的无质量标量场的最小耦合。
\end{definition}

即使在 BPS 极限下势 $V(\phi)$ 为零，规范动能与标量动能均为零、且 $\tr\langle\phi\rangle^2=v^2/2$ 的构型仍是理论的有效真空态。Bogomol'nyi 界在
\[
E_i=\sin(\theta)D_i\phi
\qquad\text{与}\qquad
B_i=\cos(\theta)D_i\phi,\qquad\text{其中 } \tan(\theta)=q_{\mathrm{el}}/q_{\mathrm{mag}}
\]
时饱和。

\begin{definition}[19.12]
这些方程称为 BPS 方程，而使 Bogomol'nyi 界饱和的构型称为 BPS 态。
\end{definition}

% [原文 p.275]

关于 BPS 极限下 Georgi--Glashow 模型的直接分析，见 [339, 第 1.3 节]。

\section{Witten 效应}

$\Theta$-项
\[
\frac{\Theta}{16\pi^2}\int \dd^4x\,F_{\mu\nu}(\star F)^{\mu\nu}
=\frac{\Theta}{8\pi^2}\int F\wedge F
\]
对场强是二次的，因而是相关的，所以可以在低能下出现。考虑把它加到式 (19.1) 中，于是完整作用量成为
\[
S_M=\int \dd^4x\left(-\frac{1}{4e^2}F_{\mu\nu}F^{\mu\nu}\right)
+\frac{\Theta}{8\pi^2}\int F\wedge F .
\]
回忆 $F$ 是复线丛 $E$ 上酉联络的曲率，有
\[
\frac{1}{8\pi^2}\int F\wedge F=\frac{1}{2}\int c_1(E)^2\in\frac{1}{2}\Z,
\]
其中 $c_1(E)$ 是 $E$ 的第一 Chern 类。

\begin{remark}[19.13]
最后一式表明 $\Theta$ 是 $4\pi$-周期的。然而，若时空是自旋流形，$c_1(E)^2$ 的积分实际上是偶数。我们将关注自旋时空，其中 $\Theta$ 事实上是 $2\pi$-周期的。
\end{remark}

\begin{remark}[19.14]
离散 Poincaré 变换 $CP$ 把 $\Theta$ 映为 $-\Theta$，所以除 $\Theta=0$ 或 $\Theta=\pi$ 外 $CP$（等价地，时间反演 $T$）被破坏。
\end{remark}

当 $\Theta\ne 0$ 时，Maxwell 方程修改为
\begin{equation}
\dd F=0, \tag{19.10}
\end{equation}
\begin{equation}
\dd\left(\frac{2\pi}{e^2}\star F+\frac{\Theta}{2\pi}F\right)=0. \tag{19.11}
\end{equation}
因此，进入电 Gauss 定律的不是电场 $\mathbf{E}$，而是
\[
\frac{2\pi}{e^2}\mathbf{E}+\frac{\Theta}{2\pi}\mathbf{B}.
\]
电荷仍然量子化，并且对任意在固定时刻包围电荷为 $ne$（其中 $n\in\Z$）的局部算子的 2-球面 $S^2$，有
\[
\int_{S^2}\left(\frac{2\pi}{e^2}\star F+\frac{\Theta}{2\pi}F\right)
=\int_{S^2}\dd\mathbf{s}\cdot\left(\frac{2\pi}{e^2}\mathbf{E}+\frac{\Theta}{2\pi}\mathbf{B}\right)=2\pi n.
\]
若绝热地改变 $\Theta$，可以看到为使 Gauss 定律继续成立，$\mathbf{E}$ 必须获得正比于 $\Theta\mathbf{B}$ 的贡献。具体地说，当 $\Theta\ne 0$ 时，

% [原文 p.276]

磁单极子会产生径向电通量，对应于电荷 $e\Theta/2\pi$。这一现象称为 Witten 效应。

\begin{proposition}[19.15]
在变换 $\Theta\to\Theta+2\pi$ 下，双子按如下方式变换：
\begin{equation}
T\colon (q_{\mathrm{el}},q_{\mathrm{mag}})\longmapsto(q_{\mathrm{el}}+q_{\mathrm{mag}},q_{\mathrm{mag}}), \tag{19.12}
\end{equation}
其中表示这一变换的 $T$ 不应与时间反演对称性混淆。
\end{proposition}

\begin{remark}[19.16]
注意，当 $\Theta\ne 0$ 时双子的电荷获得正比于 $q_{\mathrm{mag}}\Theta$ 的贡献，这与 DSZ 量子化条件并不矛盾。对电荷 $e$ 与电荷为 $(q_{\mathrm{el}},q_{\mathrm{mag}})$ 的双子，DSZ 条件要求 $q_{\mathrm{mag}}=2k\pi/e$，其中 $k\in\Z$。相反，对两个电荷为 $(q_{\mathrm{el}},q_{\mathrm{mag}})=(q_{\mathrm{el},i},2\pi/e)$（$i=1,2$）的双子，量子化条件是
\[
q_{\mathrm{el},1}-q_{\mathrm{el},2}\in e\Z .
\]
这表明，在含一个电荷为 $e$ 的粒子的一般 $\mathrm{U}(1)$ 规范理论中，DSZ 量子化条件只约束任意两个双子的电荷之差为 $e$ 的整数倍。
\end{remark}

由于 $CP$ 变换作用为 $CP\colon (q_{\mathrm{el}},q_{\mathrm{mag}})\to(-q_{\mathrm{el}},q_{\mathrm{mag}})$，$CP$-不变性把双子的电荷限制为要么全是整数（对 $\Theta=0$），要么全是半整数（对 $\Theta=\pi$）。一般 $\mathrm{U}(1)$ 规范理论中的双子，可以用前几段讨论的方法分析，只要在类 Georgi--Glashow 模型中加入破坏 $CP$ 的 $\Theta$-项。正如所预期的，人们发现双子的磁荷以因子 $e\Theta g_{\mathrm{mag}}/2\pi$ 贡献于其电荷 [361]。
% [Tutor Note: 原文 Remark 19.16 说"双子的电荷获得正比于 $q_{\mathrm{mag}}\Theta$ 的贡献"，而本段末又说磁荷"以因子 $e\Theta g_{\mathrm{mag}}/2\pi$ 贡献于电荷"（印作小写 $g_{\mathrm{mag}}$，与 $q_{\mathrm{mag}}$ 的记号混用）；两者是同一效应的两种写法，后者乘以 $e$ 即得电荷偏移 $e^2\Theta g_{\mathrm{mag}}/2\pi$，与 $q_{\mathrm{mag}}=2\pi n_m/e$ 相容。Witten 效应的标准表述是：$\Theta\to\Theta+2\pi$ 时单极子获得一个单位的电荷，这正是式 (19.12)。]

\subsection{带 $\Theta$ 项时的电磁对偶}

当 $\Theta\ne 0$ 时，电磁对偶交换式 (19.10) 与式 (19.11)。记
\begin{equation}
F_D=\frac{2\pi}{e^2}\star F+\frac{\Theta}{2\pi}F, \tag{19.13}
\end{equation}
则可以把 Maxwell 方程改写为
\[
\dd\left(\frac{2\pi}{e_D^2}\star F_D+\frac{\Theta_D}{2\pi}F\right)=0,
\qquad \dd F_D=0,
\]
只要
\begin{equation}
-\frac{4\pi^2}{e^2e_D^2}+\frac{\Theta\Theta_D}{4\pi^2}=-1,
\qquad
\frac{\Theta}{e_D^2}+\frac{\Theta_D}{e^2}=0. \tag{19.14}
\end{equation}
% [Tutor Note: 式 (19.14) 的两式恰好等价于 $\tau\tau_D=-1$：把 $\tau=2\pi\ii/e^2+\Theta/2\pi$ 与 $\tau_D=2\pi\ii/e_D^2+\Theta_D/2\pi$ 相乘，实部即第一式、虚部即第二式。因此 Prop. 19.17 中的 $\tau_D=S(\tau)=-1/\tau$ 与 (19.14) 是同一条件的两种写法，可据此核对符号约定。]
让我们把复化规范耦合与对偶规范耦合定义为：

% [原文 p.277]

\[
\tau=\frac{2\pi\ii}{e^2}+\frac{\Theta}{2\pi},
\qquad
\tau_D=\frac{2\pi\ii}{e_D^2}+\frac{\Theta_D}{2\pi}.
\]

\begin{proposition}[19.17]
电磁对偶对复化规范耦合的作用为
\[
\tau_D=S(\tau):=\frac{-1}{\tau},
\]
而对电荷的作用为
\[
S\colon (q_{\mathrm{el}},q_{\mathrm{mag}})\longmapsto(-q_{\mathrm{mag}},q_{\mathrm{el}}).
\]
\end{proposition}

用复化规范耦合表示，式 (19.12) 改写为
\[
T\colon \tau\longmapsto \tau+1 .
\]
变换 $S$ 与 $T$ 生成群 $\SL_2(\Z)$。中心元 $-\mathrm{Id}\in\SL_2(\Z)$ 对复化规范耦合 $\tau$ 的作用是平凡的，而它对电荷的作用则不然。因此，仅就耦合而言，相关的群是 $\PSL_2(\Z)=\SL_2(\Z)/\{\pm\mathrm{Id}\}$。

\section{非交换电磁对偶}

\subsection{非交换量子化条件}

't Hooft--Polyakov 与 Julia--Zee 的构造可以推广到非交换 UV 规范群 $G_0$ 自发破缺到非交换子群 $G$ 的规范理论。在这一设定下，$G_0$ Yang--Mills 理论中存在经典孤子，对应于与未破缺子群 $G$ 相联系的（非交换）磁单极子与双子。这些有限能量场构型由 Goddard、Nyuts 与 Olive 在 [165] 中研究并分类。

具体地说，设 $G_0$ 单、紧、连通且单连通，考虑带某个表示 $R$ 中标量场 $\phi$ 的 $G_0$ Yang--Mills 理论。我们关注经典真空模空间中 $G_0$ 自发破缺到单、紧、连通子群 $G$ 的子空间 $\mathcal{M}$。一般地，$G_0$ 在 $\mathcal{M}$ 上传递地作用，且 $\mathcal{M}\simeq G_0/G$。与式 (19.5) 一样，在任何有限能量 $G_0$-构型中，标量 $\phi$ 定义一个从无穷远球面 $S_\infty$ 到 $\mathcal{M}$ 的映射，它在拓扑上由其属于 $\pi_2(G_0/G)$ 的类刻画。

\begin{remark}[19.18]
任何单、紧、连通的 $G_0$ 都满足 $\pi_2(G_0)=0$。再假定 $G_0$ 单连通，则同伦群的正合列
\[
\cdots\to\pi_2(G_0)\to\pi_2(G_0/G)\to\pi_1(G)\to\pi_1(G_0)\to\pi_1(G_0/G)\to\cdots
\]

% [原文 p.278]

表明 $\pi_2(G_0/G)\simeq\pi_1(G)$。
\end{remark}

\begin{definition}[19.19]
设 $F_{\mu\nu}$ 是未破缺群 $G$ 的场强。$G$-磁单极子在其静止系中由下述条件定义：在空间无穷远处，$F$ 的空间分量为
\[
F_{ij}(\mathbf{r})=\frac{\epsilon_{ijk}r^k}{r^3}F(\mathbf{r})\in\mathfrak{g},
\]
其中 $r^k$ 是空间径向单位向量的第 $k$ 个分量，$r=|\mathbf{r}|$ 是径向距离。
\end{definition}

\begin{proposition}[19.20]
为使场论自洽，必须满足下述量子化条件：
\begin{equation}
\exp(4\pi\ii F(\mathbf{r}))=1\in G. \tag{19.15}
\end{equation}
\end{proposition}

与通常一样，必须把规范等价的构型等同起来。在未破缺子群 $G\subset G_0$ 中元素共轭的意义下，总可以假定 $F$ 落在 $G$ 的一个固定 Cartan 子群 $H$ 中。于是可以把式 (19.15) 的量子化条件重新表述为一个几何条件：

\begin{proposition}[19.21]
场强 $F(\mathbf{r})$ 必须落在 $G$-余特征格 $\Gamma^G_w\subset\mathfrak{h}$ 中，其中 $\mathfrak{h}$ 是 $H$ 的李代数。因此，$G$-单极子由其磁荷刻画，而磁荷取值于 $\Gamma^G_w/W$，其中 $W$ 是 $\mathfrak{g}=\operatorname{Lie}(G)$ 的 Weyl 群。
\end{proposition}

在低能规范群 $G$ 下“带电荷”的场必须按 $G$ 的有限维不可约表示变换。这些表示与 $G$ 的支配权一一对应，即与 $\Lambda^G_w/W$ 的元素一一对应，其中 $\Lambda^G_w\subset\mathfrak{h}^*$ 是 $G$-特征格。换言之，电荷由 $\Lambda^G_w/W$ 分类。最后，低能理论中的双子同时携带 $G$ 下的电荷与磁荷，其允许的荷与 $(\Lambda^G_w\times\Gamma^G_w)/W$ 一一对应。

\begin{remark}[19.22]
$\mathfrak{g}$ 的根是特殊的权，即 $\mathfrak{g}$ 的根格 $\Lambda_r$ 是 $\mathfrak{g}$ 的权格 $\Lambda_w$ 的子格。对满足 $\operatorname{Lie}(G)=\mathfrak{g}$ 的紧连通李群 $G$，$\mathfrak{g}$ 的表示未必都能提升为 $G$ 的表示。例如，$\mathfrak{su}(N)$ 的基本 $N$ 维表示不能提升为 $\mathrm{PSU}(N)$ 的表示。然而，$\mathfrak{g}$ 的伴随表示总能提升为 $G$ 的表示。事实上，$G$ 的有限维不可约表示与 $\Lambda^G_w/W$ 的元素一一对应，其中 $\Lambda^G_w$ 是 $G$-特征格，满足 $\Lambda_r\subset\Lambda^G_w\subset\Lambda_w\subset\mathfrak{h}^*$。
\end{remark}

对偶地，$\mathfrak{g}$ 的 Cartan 子代数 $\mathfrak{h}$ 中有两个基本格 $\Gamma_r\subset\Gamma_w\subset\mathfrak{h}$：余根格 $\Gamma_r$ 与余权格 $\Gamma_w$。群 $G$ 的选取定义了所谓的 $G$-余特征格 $\Gamma^G_w$，它满足 $\Gamma_r\subset\Gamma^G_w\subset\Gamma_w$。例如，当 $\mathfrak{g}=\mathfrak{su}(N)$ 时，有 $\Lambda^{\SU(N)}_w=\Lambda_w$ 与 $\Gamma^{\SU(N)}_w=\Gamma_r$，而 $\Lambda^{\mathrm{PSU}(N)}_w=\Lambda_r$ 与 $\Gamma^{\mathrm{PSU}(N)}_w=\Gamma_w$。李论基础汇集于附录 A。

% [原文 p.279]

\subsection{Montonen--Olive 对偶猜想}

电磁对偶必须交换带电荷的粒子与磁单极子，这一事实促使 Montonen 与 Olive 提出下述猜想。

\begin{conjecture}[19.23（Montonen--Olive 对偶 {{[271]}}）]
电磁对偶 $S$ 把规范群为 $G$、规范耦合为 $\tau$ 的纯 Yang--Mills 理论映为规范群为 $G^\vee$、规范耦合为 $-1/\tau$ 的纯 Yang--Mills 理论，其中 $G^\vee$ 是满足 $\operatorname{Lie}(G^\vee)=\mathfrak{g}=\operatorname{Lie}(G)$、且 $\Lambda^{G^\vee}_w=\Gamma^G_w$（等价地 $\Gamma^{G^\vee}_w=\Lambda^G_w$）的群。
\end{conjecture}

群 $G^\vee$ 称为 GNO 对偶群 [165]，或磁规范群（与之相对，$G$ 称为电规范群），或 Langlands 对偶群，因为它在几何 Langlands 对应中居中心地位。规范理论中的电磁对偶与几何 Langlands 对应之间的关系已在 [226] 中详尽探讨。

\begin{remark}[19.24]
在经典层面，Montonen--Olive 猜想是自然的，特别因为在任何由类 Georgi--Glashow 模型在 BPS 极限下得到的 Yang--Mills 理论中，所有经典态（包括 $W$-玻色子、单极子与双子）都使 BPS 界饱和。然而，与 Maxwell 理论不同，非交换 Yang--Mills 理论中的规范耦合会跑动，量子效应完全可能破坏 BPS 界。此外，由于缺乏分析强耦合 Yang--Mills 理论的工具，Montonen--Olive 猜想极难检验：由于 $\tau_D=-1/\tau$，对偶的一侧必然是强耦合的。稍后我们将看到，借助超对称情况会显著改善：人们可以为 Montonen--Olive 猜想提供非平凡的证据，虽然不是在纯 Yang--Mills 理论中，而是在 $\mathcal{N}=4$ 超 Yang--Mills 理论中。
\end{remark}

\subsection{规范理论的整体变体}

规范理论的规范群 $G$ 并不由物质内容决定；相反，正是规范群的选取决定了带电荷物质场可以按哪些表示变换。例如，由于伴随表示既是 $\SU(N)$ 也是 $\mathrm{PSU}(N)$ 的真正表示，$G$ 的伴随表示中的场可以同时出现在 $\SU(N)$ 与 $\mathrm{PSU}(N)$ 规范理论中。然而，基本场只能存在于 $\SU(N)$ 规范理论中，因为 $\SU(N)$ 的基本表示不是 $\mathrm{PSU}(N)$ 的真正表示。此外，即使拉格朗日量的物质内容与共享同一李代数的两个不同规范群 $G$ 与 $G'$——例如 $G=\SU(N)$ 与 $G'=\mathrm{PSU}(N)$——相容，仍可以区分对应于 $G$ 的规范理论与对应于 $G'$ 的规范理论 [7]。

这来自规范理论具有线算子这一事实，特别是 Wilson 线与 't Hooft 线。关于规范理论中 Wilson 线与 't Hooft 线的介绍，我们参阅 [346]。

% [原文 p.280]

\begin{proposition}[19.25]
Wilson 线由 $G$ 的表示标记，等价地由 $\Lambda^G_w/W$ 标记。反之，'t Hooft 线由 $\Gamma^G_w/W$ 标记。
\end{proposition}

注意这特别意味着，纯 $\SU(N)$ YM 与纯 $\mathrm{PSU}(N)$ YM 理论是不同的，因为它们没有相同的线算子内容。

更一般地，规范理论容许双子的线算子，也称为 Wilson--'t Hooft 线，其荷在 $(\Gamma^G_w\times\Lambda^G_w)/W$ 中 [225]。线算子与局部算子之间、以及两个线算子之间（称满足此条件的任意两条线是相互局部的）都有 Dirac--Schwinger--Zwanziger（DSZ）量子化条件的一个版本（满足它的线相对于动力学场是局部的），以保证理论的自洽性。$G$-规范理论的 Wilson 线与 't Hooft 线总是相互局部的。

\begin{definition}[19.26]（{[7]}）
$G$-规范理论的定义包含其线算子的集合，该集合要求是局部且相互局部的双子线的极大族（包括 Wilson 线与 't Hooft 线）。仅在线算子内容上不同的理论称为整体变体（global variants）。
\end{definition}

尽管理论中的 Wilson 线与 't Hooft 线如上所述完全由规范群 $G$ 的选取决定，一般说来，可以用双子线以不止一种方式把它们补全为极大的一族相互局部线。例如，纯 $\mathrm{PSU}(N)$ YM 理论有 $N$ 个不同的整体变体，记作 $\mathrm{PSU}(N)_0,\dots,\mathrm{PSU}(N)_{N-1}$。

线算子的荷受 $S$ 与 $T$ 作用，其中 $T$ 对应于 $\Theta$-角平移 $2\pi$。一般地，$S$ 与 $T$ 会置换整体变体。例如，在规范群非单连通的规范理论中，由于可能存在的分数瞬子，$\Theta$ 未必是 $2\pi$-周期的。因此，$T$ 把原来的整体变体映为另一个。

\section{超对称规范理论中的电磁对偶}

\subsection{Witten--Olive--BPS 态}

考虑由式 (18.3) 形式的拉格朗日量定义的 $4d$ $\mathcal{N}=2$ 渐近自由或共形规范理论。在其库仑分支的一般点处，低能有效理论是一个交换的 $4d$ $\mathcal{N}=2$ 规范理论。微观理论的有限能量规范构型在低能有效理论中给出单极子与双子。事实上，拉格朗日型 $4d$ $\mathcal{N}=2$ 规范理论是 Georgi--Glashow 理论（带费米子）在 BPS 极限下的特例。此外，BPS 界在经典极限下饱和。

有两条独立的守恒超流，这与理论具有 $\mathcal{N}=2$ 超对称相对应。把这些超流在空间切片上积分，得到守恒超荷 $Q^I_\alpha$、$\bar Q^I_{\dot\alpha}$（$I=1,2$），它们

% [原文 p.281]

作为 $4d$ $\mathcal{N}=2$ 超庞加莱代数中的奇生成元出现，并如下定义代数的中心荷 $\mathcal{Z}$：
\[
\{Q^I_\alpha,Q^J_\beta\}=2\sqrt{2}\epsilon_{\alpha\beta}\epsilon^{IJ}\mathcal{Z},
\]
这与式 (17.1) 一致。超流 $Q^I_\alpha$ 与 $\bar Q^I_{\dot\alpha}$ 的表达式通常涉及 UV Yang--Mills 场强 $F$。

回忆式 (19.8)，UV 理论孤子规范构型的电磁荷是 $F$ 在空间无穷远处的积分。在奠基性文章 [369] 中，Witten 与 Olive 证明：$4d$ $\mathcal{N}=2$ 超庞加莱代数的中心荷 $\mathcal{Z}$ 可以写成低能荷的线性组合。此外，这类理论中的 Bogomol'nyi 界读作 $m\ge\sqrt{2}|\mathcal{Z}|$，意味着它恰好与由超代数导出的式 (17.3) 的 SUSY 界一致。这解释了为何扩展超庞加莱代数的短表示被称为 BPS 态：

\begin{proposition}[19.27]（{[369]}）
在具有扩展超对称的量子场论中，BPS 态按超对称代数的短表示变换。因此，BPS 态的两个定义——一方面作为使 Bogomol'nyi 界饱和的态，另一方面作为 $4d$ $\mathcal{N}=2$ 超庞加莱代数的短表示——是一致的。
\end{proposition}

经典单极子与双子背景下的费米零模生成 $\mathcal{N}=2$ BPS 超多重态，其标量分量就是与这些孤子相应的复标量，例如见 [292] 与 [339, 第 1 章]。BPS 态与短超多重态的这一等同蕴含下述重要性质：

\begin{proposition}[19.28]
在 $4d$ $\mathcal{N}=2$ 规范理论中，BPS 界不能被量子效应破坏。因此，经典地使 Bogomol'nyi 界饱和的任何态在量子理论中也是 BPS 的。
\end{proposition}

\begin{proof}
这来自长超多重态有 16 个自由度、而短超多重态只有 8 个这一事实。
\end{proof}
% [Tutor Note: Prop. 19.28 的证明只用到"长表示 16 个自由度、短表示 8 个"这一计数，与导言章（"动机与历史回顾"，p.xxxiii 附近）"短超多重态的维数是长超多重态的一半"是同一论断。正是这一表示论保护使 BPS 态在量子理论中依然稳定，也是 $S$ 对偶能在 $\mathcal{N}=2$ 理论中成立的基础（另见 19.4.4）。]

\begin{example}[19.29]
纯 $\mathcal{N}=2$ $\SU(2)$ SYM 是渐近自由的。在其库仑分支 $\mathcal{C}$ 上远离原点的一般点处，UV 理论在自发破缺到 $\mathrm{U}(1)$ 之前从不会变强耦合。经典地，'t Hooft--Polyakov--Julia--Zee 构造给出电荷为 $(q_{\mathrm{el}},q_{\mathrm{mag}})=\pm(0,1)$ 的单极子以及电荷为 $(q_{\mathrm{el}},q_{\mathrm{mag}})=\pm(2k,1)$（任意 $k\in\Z$）的双子。因此，库仑分支上这样的点处的低能理论显示有质量的 BPS $\mathcal{N}=2$ 超多重态，其标量分量是电荷为 $(q_{\mathrm{el}},q_{\mathrm{mag}})=\pm(0,1)$ 的单极子与电荷为 $(q_{\mathrm{el}},q_{\mathrm{mag}})=\pm(2k,1)$（任意 $k\in\Z$）的双子。
\end{example}

因此，考虑 $4d$ $\mathcal{N}=2$ 规范理论（而非更一般的规范理论）解决了与 Montonen--Olive 对偶的量子解释有关的两个问题之一，即量子效应可能破坏 BPS 界这一顾虑。

% [原文 p.282]

\subsection{$4d$ $\mathcal{N}=4$ SYM 理论的 $S$ 对偶}

与纯 Yang--Mills 理论中 Montonen--Olive 对偶的量子解释有关的第二个问题是规范耦合的跑动。这可以通过改而考虑标度不变的理论来处理。四维 $\mathcal{N}=4$ 超 Yang--Mills（SYM）理论是 $4d$ $\mathcal{N}=2$ 理论的特殊情形，具体说是带伴随表示中无质量超多重态的 $4d$ $\mathcal{N}=2$ 理论。它们是超共形的，从而是标度不变的，这特别意味着规范耦合不跑动。因此，把 Montonen--Olive 猜想表述为规范群为 $G$、耦合为 $\tau$ 的 $\mathcal{N}=4$ SYM 与对偶规范群为 $G^\vee$、耦合为 $-1/\tau$ 的 $\mathcal{N}=4$ SYM 之间的对偶，是有意义的。事实上，这些理论存在不同的整体变体，这要求对该猜想作精细化的修正，稍后将予讨论。假定 Montonen--Olive 对偶成立，则两个对偶理论应被视为同一个底层量子场论的不同但等价的描述，而不是两个彼此分离的 QFT。

在 $\mathcal{N}=4$ $\SU(2)$ SYM 的库仑分支上（它是带无质量伴随超多重态的 $\mathcal{N}=2$ $\SU(2)$ SYM），低能态——电子、单极子与双子——都是 $4d$ $\mathcal{N}=4$ 超庞加莱代数的四分之一 BPS 超多重态。因此，$S$ 对偶可以把电子映为单极子、反之亦然。于是，这个理论是 Montonen--Olive 对偶实现的一个有力候选者，这一观念最初在 [292] 中提出。随后的研究，特别是 [324] 与 [348]，进一步证实了 Montonen--Olive 对偶在 $4d$ $\mathcal{N}=4$ SYM 理论中的实现，使理论界对这一现象形成了牢固的共识。

电磁对偶 $S$ 与非交换变换 $T$ 生成对偶群 $\SL_2(\Z)$\footnote{或者更确切地说，其 metaplectic 扩张 $\mathrm{Mp}_2(\Z)$ [294]。}。它对规范耦合的作用为
\[
\tau\longmapsto\frac{a\tau+b}{c\tau+d},\qquad
\text{其中 } \begin{pmatrix} a & b\\ c & d\end{pmatrix}\in\SL_2(\Z),
\]
于是对 $\tau$ 的作用归结为 $\PSL_2(\Z)$ 的作用，因为 $\{-\mathrm{Id}\}<\SL_2(\Z)$ 的作用平凡。但 $\{-\mathrm{Id}\}$ 对电荷的作用非平凡。

\begin{remark}[19.30]
$4d$ $\mathcal{N}=4$ 超 Yang--Mills（SYM）理论的 $S$ 对偶意味着，低能有效理论对一切互素的电磁荷 $(p,q)$ 都含有 BPS 态。这一论断已在 [323] 中得到直接支持。
\end{remark}

规范群为 $G$ 的规范理论存在不同的整体变体，这意味着 $S$ 对偶的陈述必须精细化。例如，规范群为 $\mathrm{SO}(3)$ 的 $\mathcal{N}=4$ SYM 有两个整体变体，记作 $\mathrm{SO}(3)_+$ 与 $\mathrm{SO}(3)_-$。$\mathcal{N}=4$ $\SU(2)$ SYM 的 $\Theta$-角是 $2\pi$-周期的，因此这个理论是 $T$-不变的。相反，$T$ 交换整体变体 $\mathrm{SO}(3)_+$ 与 $\mathrm{SO}(3)_-$。另一

% [原文 p.283]

方面，研究 $S$ 如何作用于线算子表明，它把 $\SU(2)$ 理论映为 $\mathrm{SO}(3)_+$，而 $\mathrm{SO}(3)_-$ 是 $S$-自对偶的。因此，在 $\tau=\ii$ 处，$S$ 对偶增强为 $\mathrm{SO}(3)_-$ 的对称性。$\mathfrak{su}(2)$ $\mathcal{N}=4$ SYM 理论的对偶图见图 19.1。

\Fig[0.62\linewidth]{figures/fig_19_1.png}{19.1}{$\mathfrak{su}(2)$ $\mathcal{N}=4$ SYM 理论的对偶图。}
% [Tutor Note: 读图 19.1 的方法：实线箭头 $S$ 把 $\SU(2)$ 理论映到 $\mathrm{SO}(3)_+$，而 $\mathrm{SO}(3)_-$ 上的 $S$ 箭头是自环（$S$-自对偶）；虚线箭头 $T$（即 $\Theta\to\Theta+2\pi$）交换 $\mathrm{SO}(3)_\pm$，同时是 $\SU(2)$ 上的自环（因为 $\mathcal{N}=4$ $\SU(2)$ 的 $\Theta$-角以 $2\pi$ 为周期，故该理论 $T$-不变）。自对偶点 $\tau=\ii$ 处 $S$ 成为 $\mathrm{SO}(3)_-$ 的对称性，这正是 Remark 19.31 中不可逆"对偶缺陷"出现的机制。另：正文脚注指出更精确的对偶群是 $\SL_2(\Z)$ 的 metaplectic 扩张 $\mathrm{Mp}_2(\Z)$，因为在自对偶点处 $\SL_2(\Z)$ 的作用带有额外相位。]

\begin{remark}[19.31]
在耦合 $\tau$ 的自对偶值处，把 $S$ 与一形式对称性的规范（它同样置换整体变体）结合起来，会得到不可逆的 Krammers--Wannier 型对偶缺陷 [220, 88, 221]。
\end{remark}

\subsection{$4d$ $\mathcal{N}=2$ SYM 理论的 $S$ 对偶}

为在与纯 Yang--Mills 理论更接近的语境中研究 $S$ 对偶，可以破坏一半超对称而改而考虑 $4d$ $\mathcal{N}=2$ 理论。在此语境中，BPS 态与超庞加莱代数的小表示之间的关系仍然成立，然而并非所有 $4d$ $\mathcal{N}=2$ 理论都是标度不变的。

另一个问题会出现，例如在纯 $\mathcal{N}=2$ $\SU(2)$ SYM 中。在这个理论中，沿库仑分支在低能下存在的唯一带电荷粒子是 $W$-玻色子，它们属于 $1/2$-BPS 矢量超多重态。相反，磁单极子与双子落入 $1/2$-BPS 超多重态。由于 $S$ 对偶预期把带电荷粒子与其带磁荷的对应物联系起来，这一不匹配构成一个挑战。这表明 $\mathcal{N}=2$ SYM 理论中的 $S$ 对偶（如果存在）比 Montonen--Olive 猜想所提示的更为精细。纯 $\mathcal{N}=2$ $\SU(2)$ SYM 的 Seiberg--Witten 分析（将在 20.1 节研究）表明，一个精细化的 $S$ 对偶概念确实存在。这一精细对偶涉及穿墙的概念，而就 $4d$ $\mathcal{N}=4$ SYM 的半 BPS 态而言，这一现象是不出现的。

在深入 Seiberg--Witten 理论之前，让我们讨论 $4d$ $\mathcal{N}=2$ 规范理论的库仑分支 $\mathcal{C}$ 的特殊凯勒几何与该分支上低能有效交换理论的电磁对偶之间的相互作用。回忆命题 18.25：任何秩为 $r$ 的 $\mathcal{N}=2$ 理论的库仑分支都是复维数为 $r$ 的特殊凯勒流形，这意味着凯勒势可以写成定义 18.5 中的形式：

% [原文 p.284]

\[
K(a)=\operatorname{Im}\left(\frac{\partial\mathcal{F}}{\partial a_i}\bar a_i\right),
\qquad i=1,\dots,r,
\]
其中 $a_1,\dots,a_n$ 是 $\mathcal{C}$ 上的局部全纯坐标，而 $\mathcal{F}$ 是 $a_i$ 的局部全纯函数，称为预势（prepotential）。低能交换规范耦合 $\tau^{ij}$ 如式 (18.7) 那样写成
\begin{equation}
\tau^{ij}(a)=\tau^{ji}(a)=\frac{\partial^2\mathcal{F}}{\partial a_i\partial a_j}, \tag{19.16}
\end{equation}
而特殊凯勒度量（如命题 18.6）为
\begin{equation}
\dd s^2=\operatorname{Im}\tau^{ij}(a)\,\dd a_i\,\dd\bar a_j. \tag{19.17}
\end{equation}
对任意 $i=1,\dots,r$，令
\begin{equation}
a^D_i=\frac{\partial\mathcal{F}}{\partial a_i}. \tag{19.18}
\end{equation}
用 $a_i$ 与 $a^D_i$ 表示，势 $K$ 改写为
\[
K=2\ii\left(a_i\bar a^D_i-\bar a_i a^D_i\right),
\]
而特殊凯勒度量成为
\[
\dd s^2=\operatorname{Im}\dd a^D_i\,\dd\bar a_i
=-\frac{\ii}{2}\left(\dd a^D_i\,\dd\bar a_i-\dd\bar a^D_i\,\dd a_i\right).
\]
这些表达式在 $a_i$ 与 $a^D_i$ 中对称。因此，把每个 $a_i$ 换成 $a^D_i$，将得到与式 (19.17) 类似的特殊凯勒度量，只是把不同的调和函数替换到 $\operatorname{Im}\tau(a)$ 的位置。

\begin{definition}[19.32]
变量 $a_i$ 与 $a^D_i$ 称为特殊凯勒流形 $\mathcal{C}$ 上的特殊坐标。
\end{definition}
% [Tutor Note: 自查要点：(19.16) 的对称性 $\tau^{ij}=\tau^{ji}$ 来自 $\mathcal{F}$ 的二阶偏导可交换；(19.17) 与命题 18.6 的形式一致，系数 $\operatorname{Im}\tau^{ij}$ 就是凯勒形式的分量。另：原文此处写"$a_1,\dots,a_n$ 是局部全纯坐标"，与上下文 $i=1,\dots,r$ 不符（应为 $a_1,\dots,a_r$），照印并在此指出。]

库仑分支上低能有效作用量的玻色部分为
\[
\mathcal{L}_B=\frac{\operatorname{Im}\tau^{ij}}{4\pi}\partial_\mu\bar a_i\partial^\mu a_j
+\frac{\operatorname{Im}\tau^{ij}}{8\pi}F_{\mu\nu,i}F^{\mu\nu}_j
+\frac{\operatorname{Re}\tau^{ij}}{8\pi}F_{\mu\nu,i}(\star F)^{\mu\nu}_j.
\]
推广式 (19.13)，令
\[
F^D_{i\mu\nu}=\operatorname{Im}\tau^{ij}F_{\mu\nu,j}+\operatorname{Re}\tau^{ij}(\star F)_{\mu\nu,j},
\qquad\text{以及}\qquad
\tau_{Dij}=(-\tau^{-1})_{ij}.
\]
于是可以把 $\mathcal{L}_B$ 改写为
\begin{equation}
\mathcal{L}_B=\frac{\operatorname{Im}\tau_{Dij}}{4\pi}\partial_\mu\bar a^D_i\partial^\mu a^D_j
+\frac{\operatorname{Im}\tau_{Dij}}{8\pi}F^i_{D\mu\nu}F^{\mu\nu j}_D
+\frac{\operatorname{Re}\tau_{Dij}}{8\pi}F^i_{D\mu\nu}(\star F)^{\mu\nu j}_D. \tag{19.19}
\end{equation}
现在，由于

% [原文 p.285]

\begin{align*}
\frac{\operatorname{Im}\tau^{kl}}{4\pi}\partial_\mu\bar a_k\partial^\mu a_l
&=\frac{-1}{8\pi\ii}\left(\bar\tau^{kl}-\tau^{kl}\right)\partial_\mu\bar a_k\partial^\mu a_l\\
&=\frac{-1}{8\pi\ii}\left(\bar\tau^{ki}\delta^l_i-\delta^k_j\tau^{jl}\right)\partial_\mu\bar a_k\partial^\mu a_l\\
&=\frac{1}{8\pi\ii}\left(\tau_{Dij}-\bar\tau_{Dij}\right)\bar\tau^{ki}\tau^{jl}\partial_\mu\bar a_k\partial^\mu a_l\\
&=\frac{1}{4\pi}\operatorname{Im}\tau_{Dij}\bar\tau^{ki}\partial_\mu\bar a_k\tau^{jl}\partial^\mu a_l\\
&=\frac{1}{4\pi}\operatorname{Im}\tau_{Dij}\partial_\mu\bar a^D_i\partial^\mu a^D_j,
\end{align*}
其中我们用到了式 (19.16)、(19.18) 与 (19.19)，便得到
\[
\mathcal{L}_B=\frac{\operatorname{Im}\tau_{Dij}}{4\pi}\partial_\mu\bar a^D_i\partial^\mu a^D_j
+\frac{\operatorname{Im}\tau_{Dij}}{8\pi}F^i_{D\mu\nu}F^{\mu\nu j}_D
+\frac{\operatorname{Re}\tau_{Dij}}{8\pi}F^i_{D\mu\nu}(\star F)^{\mu\nu j}_D.
\]
于是：

\begin{proposition}[19.33]
标量 $a^D_i$ 与场强为 $F^D_{i\mu\nu}$ 的磁对偶规范场 $A^D_{i\mu}$ 处于同一个 $\mathcal{N}=2$ 矢量多重态中。换言之，$a^D_i$ 在 $\mathcal{N}=2$ 电磁对偶的意义下与 $a_i$ 对偶 [313, 339]。
\end{proposition}

使式 (19.17) 的度量的一般形式不变的 $a_i$ 与 $a^D_i$ 上的变换群在 [313, 第 3 节] 中导出。设 $X\simeq\C^{2r}$ 带坐标 $a_i,a^D_i$（$i=1,\dots,r$），并配备辛形式 $\omega$ 与全纯 2-形式 $\omega_h$：
\[
\omega=\frac{\ii}{2}\left(\dd a_i\wedge \dd\bar a^D_i-\dd a^D_i\wedge \dd\bar a_i\right),
\qquad
\omega_h=\dd a_i\wedge \dd a^D_i.
\]
再设 $u_1,\dots,u_r$ 是 $\mathcal{C}$ 上 $r$ 个任意局部全纯坐标。于是局部地，全纯映射 $f\colon\mathcal{C}\to X$ 由 $2r$ 个函数 $a_i(u)$ 与 $a^D_i(u)$ 描述。固定这样的 $f$，并给 $\mathcal{C}$ 配备凯勒形式为 $f^*(\omega)$ 的度量：
\[
\dd s^2=\operatorname{Im}\frac{\partial a^D_i}{\partial u_k}\frac{\partial\bar a_i}{\partial\bar u_l}\,\dd u_k\,\dd\bar u_l.
\]
此外要求 $f^*(\omega_h)=0$，使得当对一切 $i$ 有 $u_i=a_i$ 时，所得度量具有式 (19.17) 的形式。写 $v=(a_1,\dots,a_r,a^D_1,\dots,a^D_r)^T$，则 $\mathcal{C}$ 上的度量在一切变换
\begin{equation}
v\longmapsto Mv+c \tag{19.20}
\end{equation}
下保持其特殊形式，其中 $M\in\Sp_{2r}(\R)$，而 $c$ 是常 $2r$-向量。

% [原文 p.286]

\subsection{$\mathcal{N}=2$ 中心荷、BPS 态、稳定性与穿墙}

在规范群为秩 $r$ 的 $G$ 的纯 $4d$ $\mathcal{N}=2$ SYM 中，库仑分支 $\mathcal{C}$ 上一般点处的低能态由其电磁荷 $(n^i_e,n_{m,i})$ 刻画，其中 $i=1,\dots,r$。中心荷 $\mathcal{Z}$ 在重整化下不变，由下式给出 [369, 313, 314]：
\[
\mathcal{Z}=n^i_e a_i+n_{m,i}a^D_i,
\]
其中 $(a_i,a^D_i)$ 是库仑分支 $\mathcal{C}$ 上的特殊坐标。若 UV 理论还含有物质超多重态，低能态还可以携带 $\mathrm{U}(1)$ 味荷。这些味荷对中心荷 $\mathcal{Z}$ 贡献额外项，参看式 (20.11) 与 (20.17)。

\begin{proposition}[19.34]
BPS 态一般受保护而不衰变。
\end{proposition}

\begin{proof}
考虑中心荷为 $\mathcal{Z}_1+\mathcal{Z}_2$ 的 BPS 态；其质量为 $m=\sqrt{2}|\mathcal{Z}_1+\mathcal{Z}_2|$。一般地，向量 $\mathcal{Z}_1$ 与 $\mathcal{Z}_2$ 在复平面上不共线，因此
\begin{equation}
m_1+m_2\ge\sqrt{2}\left(|\mathcal{Z}_1|+|\mathcal{Z}_2|\right)>\sqrt{2}|\mathcal{Z}_1+\mathcal{Z}_2|=m, \tag{19.21}
\end{equation}
于是该态衰变为中心荷分别为 $\mathcal{Z}_1$ 与 $\mathcal{Z}_2$、质量分别为 $m_1$ 与 $m_2$ 的组分，在能量上是不利的。
\end{proof}

\begin{remark}[19.35]
当 $\mathcal{Z}_1$ 与 $\mathcal{Z}_2$ 共线时，式 (19.21) 的严格不等式变成等式。此时，中心荷为 $\mathcal{Z}_1+\mathcal{Z}_2$ 的 BPS 态相对于衰变为中心荷为 $\mathcal{Z}_1$ 与 $\mathcal{Z}_2$ 的 BPS 组分而言，只是边际稳定的（marginally stable）。
\end{remark}

\begin{definition}[19.36]
$\mathcal{C}$ 中由要求 $\mathcal{Z}_1$ 与 $\mathcal{Z}_2$ 共线所定义的余维 1 子流形称为边际稳定墙（walls of marginal stability）。
\end{definition}

远离边际稳定墙时，中心荷为 $\mathcal{Z}_1+\mathcal{Z}_2$ 的 BPS 态受保护而不衰变。然而，越接近 $\mathcal{Z}_1$ 与 $\mathcal{Z}_2$ 共线的墙，能隙 $|\mathcal{Z}_1|+|\mathcal{Z}_2|-|\mathcal{Z}_1-\mathcal{Z}_2|$ 就越小。在墙上能隙完全闭合，BPS 态可以衰变。在墙的另一侧，BPS 态会重新稳定，然而它可能已不再存在于模空间的这一区域中。

\begin{definition}[19.37]
这一现象称为穿墙（wall-crossing）。
\end{definition}
% [Tutor Note: 式 (19.21) 中的 $m=\sqrt{2}|\mathcal{Z}_1+\mathcal{Z}_2|$ 是 $\mathcal{N}=2$ 的 BPS 质量公式；决定衰变是否可能的是能隙，而本段把能隙印作 $|\mathcal{Z}_1|+|\mathcal{Z}_2|-|\mathcal{Z}_1-\mathcal{Z}_2|$——按三角不等式与上下文，它应为 $|\mathcal{Z}_1|+|\mathcal{Z}_2|-|\mathcal{Z}_1+\mathcal{Z}_2|$，这样在 $\mathcal{Z}_1,\mathcal{Z}_2$ 共线时能隙才恰好闭合（与 (19.21) 取等号同步）。译文照印原文记号，并在此给出正确形式。另：原文的中心荷记号 $Z$ 与 $\mathcal{Z}$ 混用，译文统一用 $\mathcal{Z}$。]
